\section*{Bab \ref{ch:g5:speed-distance-time} --- Kelajuan: Jarak dan Waktu}

\begin{solution}{exo:g5:speed-distance-time:1}
Jarak yang sama --- waktu yang lebih sedikit menang; waktu yang sama
--- jarak yang lebih jauh menang. ``Aku berlari lebih jauh'' tidak
menetapkan satu pun: dengan waktu yang tidak sama, jarak yang lebih
jauh tidak membuktikan apa-apa tentang menjadi cepat.
\end{solution}

\begin{solution}{exo:g5:speed-distance-time:2}
Bo --- jarak yang sama (satu kolam), waktu yang lebih sedikit
($25 < 30$ sekon): aturan pertama.
\end{solution}

\begin{solution}{exo:g5:speed-distance-time:3}
Skuternya --- waktu yang sama (satu \omterm{ex:g2:measuring-time:clock}{jam}), jarak yang lebih jauh
($45 > 20$ kilometer): aturan kedua.
\end{solution}

\begin{solution}{exo:g5:speed-distance-time:4}
Pada setiap \omterm{ex:g2:measuring-time:clock}{jam} ia menempuh \qty{200}{km}. Dalam $2$ \omterm{ex:g2:measuring-time:clock}{jam}: $200 \times 2
= \qty{400}{km}$. Dalam setengah \omterm{ex:g2:measuring-time:clock}{jam}: separuh dari $200$, jadi \qty{100}{km}.
\end{solution}

\begin{solution}{exo:g5:speed-distance-time:5}
Siput ($5$ \omterm{def:g2:measuring-length:units}{meter} tiap \omterm{ex:g2:measuring-time:clock}{jam}), pejalan kaki ($5$ kilometer), bus kota
($30$), merpati ($80$ kilometer tiap \omterm{ex:g2:measuring-time:clock}{jam}).
\end{solution}

\begin{solution}{exo:g5:speed-distance-time:6}
Satu \omterm{ex:g2:measuring-time:clock}{jam} berisi $30$ jatah $2$ menit: $30 \times 100 =
\qty{3000}{m}$ --- \omterm{def:g5:speed-distance-time:speed}{kelajuan} sebesar $3$ kilometer tiap \omterm{ex:g2:measuring-time:clock}{jam}.
\end{solution}

\begin{solution}{exo:g5:speed-distance-time:7}
Mobilnya: $80 \times 3 = \qty{240}{km}$. Pesepedanya: $15 \times 3 =
\qty{45}{km}$.
\end{solution}

\begin{solution}{exo:g5:speed-distance-time:8}
Garis pengembara yang lebih cepat memanjat lebih curam. Sesudah $2$
\omterm{ex:g2:measuring-time:clock}{jam} garis pesepedanya berdiri di \qty{30}{km}.
\end{solution}

\begin{solution}{exo:g5:speed-distance-time:9}
\omterm{ex:g2:measuring-time:clock}{Jam} demi \omterm{ex:g2:measuring-time:clock}{jam}: $30$, $60$, $90$ --- sesudah $3$ \omterm{ex:g2:measuring-time:clock}{jam} ferinya sudah
menempuh tepat \qty{90}{km}: penyeberangannya memakan $3$ \omterm{ex:g2:measuring-time:clock}{jam}.
\end{solution}

\begin{solution}{exo:g5:speed-distance-time:10}
Tiap \omterm{ex:g2:measuring-time:clock}{jam} Mara menutup \qty{20}{km} dan Jon \qty{10}{km}: jaraknya
menyusut \qty{30}{km} tiap \omterm{ex:g2:measuring-time:clock}{jam}. Dari \qty{60}{km}, ia mencapai nol
sesudah $2$ \omterm{ex:g2:measuring-time:clock}{jam} --- mereka bertemu, \qty{40}{km} dari kota Mara.
\end{solution}

\begin{solution}{exo:g5:speed-distance-time:11}
Sesudah satu menit berlari, kelincinya berdiri di \qty{60}{m} ---
separuh jalan --- lalu tertidur. Kura-kuranya memerlukan
$120 \div 2 = 60$ menit dan tiba pada menit $60$. Kelincinya bangun
pada menit $1 + 70 =
71$, berlari \qty{60}{m} terakhirnya dalam satu menit, dan tiba pada
menit $72$. Kura-kuranya menang dengan selisih $12$ menit --- yang
ajek mengalahkan yang cepat, sekarang dengan aritmetika.
\end{solution}
